\section*{Conclusion}
\addcontentsline{toc}{section}{\protect\numberline{}Conclusion}
\label{sec:conclusion}

Le travail présenté dans ce rapport s’inscrit dans une démarche d’ingénierie exploratoire visant à
évaluer la faisabilité et la pertinence expérimentale d’un système de séparation d’étages reposant sur
le principe de l’aérofreinage, dans le cadre d’un lanceur expérimental. L’objectif du Projet Master
Innovation n’était pas de proposer une solution optimisée ou industrialisable, mais de déterminer si un
tel principe pouvait être mis en œuvre dans des conditions réalistes et conduire à l’acquisition de
données exploitables.\\

L’étude menée a permis, dans un premier temps, de positionner le système étudié au regard de l’état de
l’art des architectures de séparation existantes, en mettant en évidence la diversité des solutions
employées et les compromis associés. Cette analyse a motivé le choix d’une approche par aérofreinage,
encore peu explorée dans un cadre expérimental étudiant, et a conduit à la définition d’un dispositif
spécifique intégrant à la fois des éléments mécaniques, électroniques et logiciels.\\

La conception du lanceur SP-02 comme vecteur expérimental dédié a constitué un élément structurant du
projet. Celui-ci a été pensé dès l’origine pour générer des conditions de vol représentatives d’un
lanceur expérimental multi-étages, tout en intégrant les moyens nécessaires à la mise en œuvre, à la
commande et à l’observation de l’expérience de séparation. L’architecture retenue permet ainsi de
dissocier clairement le rôle du lanceur, en tant que support, de l’objet principal de l’étude.\\

Les travaux de conception et de simulation ont permis de caractériser les efforts aérodynamiques
générés par les dispositifs d’aérofreinage et d’en analyser l’impact sur la dynamique du système. Les
analyses structurelles associées ont quant à elles permis de vérifier la cohérence mécanique globale
des choix retenus et d’identifier les zones critiques à surveiller. Ces résultats, bien que dépendants
des hypothèses de modélisation, constituent une base solide pour l’évaluation du principe étudié.\\

Un apport majeur du projet réside également dans la mise en place d’une chaîne de mesure embarquée
robuste et cohérente, reposant sur la plateforme APEX. L’architecture logicielle, fondée sur une
approche orientée données et temps réel, la synchronisation temporelle des mesures, ainsi que la
gestion des flux et des télécommunications, ont été pensées pour garantir l’exploitabilité des données
lors d’un événement transitoire critique. Cette chaîne constitue un outil central pour la validation
physique du système et, plus largement, un socle réutilisable pour de futurs projets expérimentaux.\\

Les résultats présentés montrent que la mise en œuvre d’un système de séparation par aérofreinage est
techniquement envisageable dans un cadre expérimental tel que celui de SP-02. Néanmoins, plusieurs
limites subsistent, notamment liées aux incertitudes de modélisation, aux conditions réelles de vol et
à la complexité inhérente à l’analyse de phénomènes fortement transitoires. Ces éléments soulignent
l’importance de la phase expérimentale pour confronter les hypothèses retenues aux comportements
observés.\\

Au-delà de la validation du principe étudié, ce projet a permis l’acquisition de compétences et de
connaissances nouvelles au sein de l’association AéroIPSA, en particulier dans le domaine des lanceurs
multi-étages et des systèmes avioniques expérimentaux. Les travaux réalisés constituent ainsi un
premier jalon, ouvrant la voie à des études complémentaires, tant sur l’optimisation du dispositif de
séparation que sur l’exploitation approfondie des données issues des essais en vol.\\

En conclusion, ce PMI a permis de démontrer la faisabilité globale de l’expérience envisagée et de
poser un cadre méthodologique et technique cohérent pour l’étude de systèmes de séparation innovants
dans un contexte expérimental. Les perspectives ouvertes par ce travail dépassent le seul cadre du
présent projet et contribuent plus largement à la structuration de futures expérimentations au sein de
l’association.

\newpage
